% Ensayo de Hipolito

En este artículo se aborda la vulnerabilidad de nuestra información personal ante la creciente presencia de IoT en nuestra vida cotidiana a través de los años. Pues cada año aumenta la cantidad de personas que tienen acceso a un dispositivo electrónico y al internet, lo cuál ocasiona que cada año, empresas como: Facebook, Twitter, Instagram, Pinterest, etc tengan acceso a mayor cantidad de datos personales de los usuarios. También aborda la relevancia e impactos que tiene la información que cada persona ofrece a monopolios digitales a cambio del servicio que prestan dichas empresas y a no ser excluido socialmente por no tener acceso a dichos servicios. Toda esta información (Big Data) significa grandes ingresos para los monopolios digitales que la almacenan, contrastando enormemente con la cantidad que estás empresas la evalúan (0,0005 dólares por persona por información general). 
La importancia de ser conscientes de toda la información que proporcionamos por un servicio y temas sociales, al igual que lo que representa para estos monopolios digitales radica exactamente en el contenido de esos datos y a quien son vendidos en masa. Si bien la información proporcionada puede llegar a facilitar o agilizar operaciones que queramos hacer en internet, también es cierto que hay información privada y delicada (objetivamente o subjetivamente) a la cual no solo tienen acceso estás plataformas, sino también empresas a las cuales estás plataformas venden tu información, eliminando la privacidad de tu información. 
El problema con la comercialización de la información personal que estos monopolios digitales llevan a cabo, no solo es la privacidad y vulnerabilidad de la misma información. También muestra que la arquitectura actual de manejo de datos prioriza más remuneración del monopolio que el control de la información del usuario dueño de esa información.


Apoyándose en cifras de instituciones como la Unión Internacional de Telecomunicaciones y el Banco Mundial, la autora argumenta cómo el crecimiento exponencial de los datos personales (Big Data), recolectados a través de la infraestructura del IoT, ha mercantilizado nuestra identidad. Nuestros datos se han convertido en el motor económico del mercado digital, generando la mayor parte de las ganancias para las empresas tecnológicas. Ante esto, el objetivo central del texto no es persuadirnos de abandonar Internet ,lo cual sería imposible en la vida cotidiana actual, sino concientizarnos sobre el valor real de nuestra información. La autora busca convencernos de que este empoderamiento es el paso fundamental para impulsar tecnologías centradas en el usuario, como las bóvedas de datos personales (Personal Data Stores) o el proyecto The HAT, mediante las cuales el individuo puede recuperar el control de su propia información. 
La temática central del artículo gira en torno a la mercantilización de la información, mientras que de forma lateral aborda el desarrollo del IoT y el manejo de la privacidad del usuario. Esta problemática exige que los estudiantes de Ciencias de la Computación entendamos a fondo estas arquitecturas. La recolección masiva de datos en el Big Data requiere Sistemas Manejadores de Bases de Datos robustos. Al momento de diseñar estas bases de datos en mi futura práctica profesional, debo aplicar el criterio técnico y ético para proteger la privacidad del usuario en la medida de lo posible, garantizando a través de un buen esquema la seguridad, el aislamiento y la independencia de los datos.


A nivel personal, es fácil caer en el error de subestimar el valor de nuestros propios datos. Por ejemplo, al analizar mi trayecto diario a la universidad, inicialmente podría pensar que mi geolocalización o mis horarios de transporte carecen de importancia para terceros. Sin embargo, analizado bajo la óptica del Big Data, este flujo constante de información revela mi patrón de vida, mi nivel socioeconómico y mis hábitos de consumo de forma detallada. Al regalar esta información por el simple uso de una aplicación móvil, confirmo la postura de la autora: los usuarios alimentamos a los monopolios digitales sin ser conscientes del poder económico que les estamos cediendo. Por lo tanto, el desarrollo de herramientas donde el usuario decida qué comparte, como los Personal Data Stores, es una evolución necesaria para la privacidad digital.
Necesitamos tener en cuenta lo que vale nuestra información privada y que tan privada es. Así podemos diseñar arquitecturas de bases de datos que se concentren aún más en el usuario o al menos utilizar las herramientas que sean posibles para tener el control de la información que proporcionamos. Aprendí que no debemos subestimar la centralización de nuestros datos personales, lo fácil que es proporcionarlos y toda la información que se pueda derivar de la que nosotros damos. Me parece que todo lo anterior deja las siguientes preguntas: ¿Qué información si estamos dispuestos a ofrecer? Mientras menos información demos, ¿Qué porcentaje del servicio se reducirá? ¿Tener un criterio ético al diseñar arquitecturas de bases de datos afectará nuestro trabajo? ¿En realidad tendremos la libertad de tener criterios éticos?

